\section*{Chapitre \ref{ch:g7:prop} --- Proportionnalité}

\begin{solution}{exo:g7:prop:1}
Premier tableau~: \omterm{def:g3:division:remainder}{quotients} $\frac{7.5}{3} = 2.5$,
$\frac{17.5}{7} = 2.5$, $\frac{27.5}{11} = 2.5$~: proportionnel,
coefficient $2.5$.

Second tableau~: $\frac62 = 3$ et $\frac{15}{5} = 3$, mais
$\frac{28}{9} \approx 3.1$~: non proportionnel.
\end{solution}

\begin{solution}{exo:g7:prop:2}
Par la règle de trois~: $5 \times p = 6 \times 8$, donc
$p = \frac{48}{5} = 9.60$ euros.
\end{solution}

\begin{solution}{exo:g7:prop:3}
Rythme~: $36 \div 3 = 12$ pages par minute. En $5$ minutes~:
$60$ pages. Pour $96$ pages~: $96 \div 12 = 8$ minutes.
\end{solution}

\begin{solution}{exo:g7:prop:4}
$30\,\%$ de $250$~: $0.3 \times 250 = 75$. \quad
$15\,\%$ de $60$~: $0.15 \times 60 = 9$. \quad
$t\,\%$ de $80$ vaut $20$~: $\frac{20}{80} = \frac{25}{100}$, donc
$t = 25$.
\end{solution}

\begin{solution}{exo:g7:prop:5}
$\frac{18}{25} = \frac{18 \times 4}{25 \times 4} = \frac{72}{100} =
72\,\%$.
\end{solution}

\begin{solution}{exo:g7:prop:6}
Chemin~: $9 \times 25\,000 = 225\,000$ cm $= 2.25$ km.
Lac~: $2$ km $= 200\,000$ cm au sol, donc
$200\,000 \div 25\,000 = 8$ cm sur la carte.
\end{solution}

\begin{solution}{exo:g7:prop:7}
$4.3$ m $= 430$ cm, et $430 \div 43 = 10$ cm.
\end{solution}

\begin{solution}{exo:g7:prop:8}
Chocolat par part~: $240 \div 8 = 30$ g. Avec $300$ g~:
$300 \div 30 = 10$ parts exactement.
\end{solution}

\begin{solution}{exo:g7:prop:9}
\emph{1.} Réduction~: $20$ euros, soit $\frac{20}{80} = 25\,\%$ du
prix initial.

\emph{2.} La hausse de $20$ euros se compare maintenant à la
\emph{nouvelle} référence $60$~: $\frac{20}{60} = \frac13 \approx
33\,\%$. Un pourcentage se rapporte toujours à un tout~; le tout a
changé, donc le pourcentage aussi.
\end{solution}

\begin{solution}{exo:g7:prop:10}
Coefficient d'après la colonne complète~: $\frac{17.5}{10} = 1.75$.
Puis $y = 4 \times 1.75 = 7$ et $x = \frac{10.5}{1.75} = 6$.
\end{solution}

\begin{solution}{exo:g7:prop:11}
\emph{1.} Après une copie~: $10 \times 0.8 = 8$ cm. Après deux~:
$8 \times 0.8 = 6.4$ cm.

\emph{2.} La seconde réduction s'applique à la copie déjà réduite, et
non à l'original~: les facteurs se multiplient,
$0.8 \times 0.8 = 0.64$, donc deux copies réduisent à $64\,\%$ de
l'original --- pas $60\,\%$.
\end{solution}

\begin{solution}{pb:g7:prop:1}
\textbf{1.} Hauteurs et ombres sont \omterm{def:g6:propdata:def}{proportionnelles}, donc par le
\omterm{def:g2:mult:def}{produit} en croix (\cref{thm:g7:prop:cross}), avec la colonne du
bâton $(1,\ 0.4)$ et celle de l'arbre $(h,\ 3.2)$~:
$0.4 \times h = 1 \times 3.2$, d'où
$h = \frac{3.2}{0.4} = 8$ m.

\textbf{2.} Tour~: $\frac{26}{0.4} = 65$ m. Ombre de l'enfant~:
$1.5 \times 0.4 = 0.6$ m (ombre $=$ hauteur $\times$ le coefficient
$0.4$ de l'instant).

\textbf{3.} À mesure que le Soleil traverse le ciel, le coefficient
qui relie les hauteurs aux ombres change~; seules des mesures prises
au même instant partagent le même coefficient.

\textbf{4.} À cet instant, le coefficient vaut $1$~: toute hauteur
est égale à son ombre. Le \omterm{def:g5:solids:def}{sommet} de la \omterm{def:g5:solids:def}{pyramide} culmine donc à
$147$ m --- la hauteur de la grande \omterm{def:g5:solids:def}{pyramide} (l'extrémité de son
ombre, mesurée depuis le point situé sous le \omterm{def:g5:solids:def}{sommet}, est tout ce
dont on a besoin).

\textbf{5.} Le \omterm{def:g3:division:remainder}{quotient} d'une colonne
$\frac{\text{ombre}}{\text{hauteur}}$ --- le coefficient de
\omterm{def:g7:prop:table}{proportionnalité} de l'instant --- est commun à tous les objets
verticaux~: c'est exactement ce qui fait du tableau des hauteurs et
des ombres un \omterm{def:g7:prop:table}{tableau de proportionnalité}
(\cref{def:g7:prop:table}).

\textbf{6.} Les \omterm{def:g3:shapes:shapes}{rayons} du Soleil arrivent \omterm{def:g6:lines:perp}{parallèles}. Sur une Terre
\emph{plate}, deux \omterm{def:g3:shapes:shapes}{rayons} \omterm{def:g6:lines:perp}{parallèles} frapperaient deux bâtons
verticaux sous le même angle~: tous deux projetteraient des ombres
\omterm{def:g6:propdata:def}{proportionnelles} --- soit aucune ombre pour les deux, soit une ombre
pour les deux. Un bâton sans ombre (Syène) et un bâton avec ombre
(Alexandrie) au même instant sont donc impossibles sur une Terre
plate~: les deux verticales pointent forcément dans des directions
différentes --- la surface est courbe.

\textbf{7.} Sur le dessin, la verticale de Syène pointe droit vers
le Soleil~; celle d'Alexandrie est inclinée de l'angle entre les
directions des deux villes, vues depuis le centre. Comme les \omterm{def:g3:shapes:shapes}{rayons}
sont \omterm{def:g6:lines:perp}{parallèles}, cette inclinaison est exactement les $7.2^\circ$
mesurés entre le \omterm{def:g3:shapes:shapes}{rayon} et le bâton à Alexandrie.

\textbf{8.} $\frac{7.2}{360} = \frac{72}{3600} = \frac{1}{50}$~: un
cinquantième de tour complet.

\textbf{9.} Les longueurs d'arc sont \omterm{def:g6:propdata:def}{proportionnelles} aux angles~:
si $7.2^\circ$ --- un cinquantième du tour --- correspondent à
$800$ km, le tour complet correspond à
\[
50 \times 800 = 40\,000 \text{ km}.
\]

\textbf{10.} \omterm{def:g4:geometry:circle}{Diamètre} $= \frac{\text{circonférence}}{\pi}
\approx \frac{40\,000}{3.14} \approx 12\,700$ km --- contre
$12\,742$ km aujourd'hui~: juste à bien moins d'un pour cent près,
avec un bâton.

\textbf{11.} $40\,000$ km $= 4\,000\,000\,000$ cm, et
$4\,000\,000\,000 \div 40\,000\,000 = 100$ cm $= 1$ m de
circonférence. \omterm{def:g4:geometry:circle}{Diamètre}~: $100 \div 3.14 \approx 32$ cm --- un beau
globe de bureau.

\textbf{12.} $4.8$ km $= 480\,000$ cm, et
$480\,000 \div 40\,000\,000 = 0.012$ cm $= 0.12$ mm. Le plus haut
\omterm{def:g5:solids:def}{sommet} des Alpes fait un dixième de millimètre sur un globe d'un
mètre de tour~: un grain de poussière. À cette \omterm{def:g7:prop:scale}{échelle}, la Terre est
plus lisse que la plupart des billes polies.

\textbf{13.} $40\,000 \div 40 = 1\,000$ jours, et
$1\,000 \div 365 \approx 2.7$~: presque trois ans de marche.

\textbf{14.} $\frac{10}{6\,370} \approx 0.0016$, soit environ
$0.2\,\%$ (plus précisément $0.16\,\%$)~: l'atmosphère est une peau
d'une finesse de souffle sur la planète.

\textbf{15.} Erreur~: $42\,000 - 40\,000 = 2\,000$ km, et
$\frac{2\,000}{40\,000} = \frac{5}{100} = 5\,\%$. En une phrase~:
avec un bâton, un puits, une route mesurée et une seule
\omterm{def:g7:prop:table}{proportionnalité}, Ératosthène a mesuré une planète à quelques pour
cent près --- les mathématiques voyagent loin avec très peu.
\end{solution}
